Atalla-Sim demonstrates that a custom cycle-level simulator can expose meaningful architectural bottlenecks before a complete RTL environment is available. In the $1024 \times 1024$ fp16 GEMM study, it showed that the naive tiled schedule is preload-limited, that blocked M/N reuse materially reduces VLS-visible traffic, and that the best measured configuration, $m8\_n32$ with scratchpad frontend queue size 4, reaches 11.22M cycles for a 2.04$\times$ speedup over the naive baseline. These results validate the simulator as a design-space exploration tool.

Although Atalla-Sim provides cycle-accurate visibility into queue pressure, overlap, and utilization, sweeping multiple parameters on the $1024 \times 1024$ GEMM still takes several hours, making broader exploration expensive and slows iteration of scheduling and memory ideas. Each new experiment currently depends on a handwritten harness that encodes staging policy, scratchpad placement, launch order, completion detection, and statistics collection in hundreds of lines of Python. This harness-centric workflow is hard to maintain, validate, and is too cumbersome to scaling to richer workloads.

The most important next step is therefore to repurpose the existing scheduler functional simulator into a cycle-accurate execution frontend for Atalla-Sim. The simulator should accept compiler-generated assembly files directly and run the same scheduled programs used by the software stack. That would remove a large amount of bespoke experiment code, align functional validation with performance evaluation, and make it practical to study more workloads than a single hand-assembled GEMM benchmark.

The current DRAM model is intentionally abstract, which is useful for early architectural exploration but limits the quality of memory-system conclusions. Hooking Ramulator into the system in place of this abstracted DRAM model would provide more realistic timing behavior and expose more informative performance metrics, including average memory access latency, read and write queue occupancy, row-buffer hit and miss behavior, bank-level parallelism, conflict-induced stalls, and sustained bandwidth under realistic request scheduling \cite{ramulator}.

With a compiler-driven execution path and higher-fidelity DRAM timing, Atalla-Sim can evolve from a useful prototype into a more scalable and more trustworthy platform for pre-silicon performance analysis.
